% Generated draft: requires mathematical and source review.
\documentclass[11pt]{article}
\usepackage[margin=25mm]{geometry}
\usepackage{fontspec}
\setmainfont{Latin Modern Roman}
\usepackage{amsmath,amssymb,mathtools,unicode-math}
\setmathfont{Latin Modern Math}
\usepackage{enumitem,needspace,xurl,hyperref}
\setlength{\parindent}{0pt}
\setlength{\parskip}{5pt}
\setlength{\emergencystretch}{3em}
\newcommand{\N}{\mathbb{N}}
\newcommand{\Z}{\mathbb{Z}}
\newcommand{\Q}{\mathbb{Q}}
\newcommand{\R}{\mathbb{R}}
\newcommand{\C}{\mathbb{C}}
\newcommand{\F}{\mathbb{F}}
\newcommand{\A}{\mathbb{A}}
\newcommand{\PP}{\mathbb{P}}
\newcommand{\E}{\mathbb{E}}
\newcommand{\eps}{\varepsilon}
\newcommand{\dd}{\,\mathrm{d}}
\begin{document}
\section*{UnsolvedMath ID: 700234}
\textbf{Source problem code:} catalogue entry
\subsection{Definitions and mathematical statement}
Let \(G\) be a finite simple group and let \(T\in\operatorname{Syl}_2(G)\) be nonabelian. For a finite group \(H\), write \(Z(H)\) for its center, \(\exp(H)\) for the least positive integer \(e\) such that \(h^e=1\) for every \(h\in H\), \(H'=[H,H]\) for its derived subgroup, and \(\Phi(H)\) for the intersection of its maximal proper subgroups. Define \(\gamma_1(H)=H\) and \(\gamma_{i+1}(H)=[\gamma_i(H),H]\). The nilpotency class of \(T\) is the integer \(n\geq 2\) satisfying \(\gamma_n(T)\neq 1\) and \(\gamma_{n+1}(T)=1\).

\textbf{(a) Sharp exponent bound.} For each integer \(n\geq 2\) that occurs as the nilpotency class of such a Sylow subgroup, determine
\[
B(n)=\sup\bigl\{\exp(S): H\text{ is a finite simple group},\ S\in\operatorname{Syl}_2(H),\ S\text{ is nonabelian of nilpotency class }n\bigr\}.
\]
Equivalently, determine the best possible upper bound for \(\exp(T)\) in terms of its nilpotency class. The supplied source states that the sharp upper bound for \(\exp(Z(T))\) is \(2^{n-1}\), and that this implies
\[
\exp(T)\leq 2^{n(n-1)}.
\]
It describes the latter bound as almost certainly too crude. These assertions and the claimed implication are retained as source claims, without independent verification.

\textbf{(d) A small family controlling fusion.} For \(U\leq G\), define
\[
N_G(U)=\{g\in G:g^{-1}Ug=U\},\qquad x^g=g^{-1}xg.
\]
For subgroups \(T_1,\ldots,T_m\leq T\), say that their normalizers together control fusion in \(T\) if the following holds. For every \(P\leq T\) and every \(g\in G\) with \(P^g\leq T\), there exist an integer \(r\geq 0\), subgroups \(P_0,\ldots,P_r\leq T\), indices \(j_1,\ldots,j_r\in\{1,\ldots,m\}\), and elements \(h_i\in N_G(T_{j_i})\) such that
\[
P_0=P,\qquad P_r=P^g,\qquad P_{i-1}^{h_i}=P_i\quad(1\leq i\leq r),
\]
and
\[
x^g=x^{h_1\cdots h_r}\qquad\text{for every }x\in P.
\]
For \(r=0\), the product is the identity. This expresses control of fusion through restrictions and compositions of conjugations, in the sense indicated by the source's reference to Alperin.

The requested task is to find a ``small number'' \(m\) of subgroups \(T_1,\ldots,T_m\) which depend only on the isomorphism class of \(T\), such that \(N_G(T_1),\ldots,N_G(T_m)\) together control fusion in \(T\) for every finite simple group \(G\) having \(T\) as a Sylow \(2\)-subgroup. The source supplies neither a numerical meaning of ``small'' nor an exact invariance requirement for the chosen family under isomorphisms of \(T\). These qualifications remain unresolved; no particular bound, characteristic-subgroup requirement, or canonical selection rule is imposed here. The number of selected subgroups is denoted by \(m\) to distinguish it from the nilpotency class \(n\) in part (a).

\textbf{Archived parts (b) and (c).} Under the same hypotheses on \(G\) and \(T\), the original additional questions were:
\begin{itemize}
\item[(b)] Can \(T\) be an internal direct product of two proper subgroups? Explicitly, can there exist \(A,B<T\) with \(T=AB\), \(A\cap B=1\), and \([A,B]=1\)?
\item[(c)] Must \(T'=\Phi(T)\) hold?
\end{itemize}
The supplied archive reports an affirmative answer to (b): Sylow \(2\)-subgroups of the alternating groups \(\mathbb{A}_{14}\) and \(\mathbb{A}_{15}\) are isomorphic to a Sylow \(2\)-subgroup of \(\mathbb{S}_4\times\mathbb{S}_8\). It also reports that Sylow \(2\)-subgroups of the groups denoted \(D_4(q)\), for odd \(q\), decompose as direct products. The precise convention for \(D_4(q)\) is not specified in the supplied record. For (c), the archive reports a negative answer, giving \(G=\operatorname{PSL}_3(q)\) with \(q\equiv 1\pmod 4\) as counterexamples; here \(q\) is a prime power, as required by the finite-field notation. These are supplied archived conclusions, not independently checked results. The supplied registry classifies the overall problem as partially solved, with (a) and (d) remaining in the main problem section.
\subsection{Short English statement}
For a nonabelian Sylow \(2\)-subgroup \(T\) of a finite simple group, determine the sharp exponent bound in terms of the nilpotency class, and find a small family of subgroups determined by the abstract group \(T\) whose normalizers control fusion. The meaning of ``small'' and the required isomorphism invariance are unspecified. The supplied archive records answers to parts (b) and (c), while the registry labels the overall problem partially solved.
\subsection{Sources}
\begin{itemize}
\item Local UnsolvedMath registry, supplied complete source record: permanent ID \(700234\), exact source problem code \texttt{catalogue entry}, title \emph{Kourovka Notebook Problem 4.24}. Proposed by D. M. Goldschmidt in 1973. Supplied upstream revision: \texttt{d371ee917ac9647fbb06cd6fc529deec6f3dda82}.
\item \emph{The Kourovka Notebook (Unsolved Problems in Group Theory)}, 4th issue (1973), Problem 4.24; the supplied record describes an update in the 21st issue (2026), identified as \texttt{arXiv:1401.0300v48}. The registry describes the Issues 1--21 collection as edited by E. I. Khukhro and V. D. Mazurov. No page number or separate bibliographic reference for Alperin's fusion terminology is supplied.
\item Supplied \emph{Archive of Solved Problems} entry for part (b), attributed to A. S. Kondratiev, \emph{Letter of October, 13, 1977}: alternating-group examples and the additional assertion concerning \(D_4(q)\) for odd \(q\).
\item Supplied archive entry for part (c), attributed to A. S. Kondratiev: counterexamples \(\operatorname{PSL}_3(q)\) with \(q\equiv 1\pmod 4\). No separate publication or date is supplied for this entry.
\item The supplied status review is dated \texttt{2026-10-16}, records evidence as \texttt{published}, and assigns the canonical status \texttt{partially\_solved}. This is supplied dated curation, not an independently verified literature assessment. Its explanatory excerpt is truncated; the fuller archived statements above come from the supplied background field.
\end{itemize}
\end{document}
